\cotefile{93}
\scanpage{47}
% archive foliation: 46
peut énoncer qu'il y a un isomorphisme canonique (fonctoriel, compatible aux extensions de la base)
\[
\deleted{\unsure{\underline{E}(}}\mathbb{D}^{*}(A)=\underline{\mathrm{Lie}}(\underline{E}(A^{*})),
\]
la structure d'extension de $\mathbb{D}^{*}(A)$ n'étant autre que la structure\margin{(sur \uline{Lie})} déduite de lastructure d'extension de $\underline{E}(A)_S=E(A)$.

Ceci posé, revenant aux conditions du théorème de déformations, on voit, par un sorite sur l'exponentielle sur lequel nous reviendrons, que la donnée d'un prolongement $\mathrm{Fil}^{1}\subset\mathbb{D}^{*}(A)_{S'}=\underline{\mathrm{Lie}}(\underline{E}(A^{*})_{S'})$ localement facteur direct prolongeant $\underline{\mathrm{Lie}}(\unsure{\underline{\underline{t}_A}}^{\inserted{\vee}})=\underline{t}_A^{\inserted{\vee}}\subset\mathbb{D}^{*}(A)_S=\underline{\mathrm{Lie}}(\underline{E}(A^{*})_{S'}\mid S)$ équivaut à la donnée d'un sous-groupe \inserted{lisse} à structure vectorielle de $\underline{E}(A^{*})_{S'}$ qui prolonge le sous-groupe $\underline{t}_A^{\inserted{\vee}}$ de $\underline{E}(A^{*})_S$. Si L est ce sous-groupe, alors il est immédiat que $\deleted{\unsure{\mathbb{D}^{*}(}}$
\[
\underline{E}(A^{*})_{\inserted{\unsure{S'}}}/L
\]
est un schéma abélien sur S', dont le schéma abélien dual est le A' cherché. (NB Il serait plus naturel ici de travailler avec $\deleted{\unsure{\mathbb{D}_{*}(A)=\mathbb{D}(A)}}$
\[
\mathbb{D}_{*}(A)=\underline{\mathrm{Lie}}\,\underline{E}(A)\xrightarrow{\substack{\mathrm{isom\ can}\\\unsure{\sim}}}\mathbb{D}^{*}(A^{*})^{\inserted{\unsure{\vee}}}
\]
plutôt qu'avec $\mathbb{D}^{*}(A)$ pour énoncer le théorème de déformations; l'équivalence des deux points de vue provient de l'accouplement parfait
\[
\mathbb{D}^{*}(A)\otimes\mathbb{D}^{*}(A^{*})\longrightarrow\underline{O}_{S_{\mathrm{\unsure{crirsnilp|crisnilp}}}}
\]
compatible avec les filtrations.)

\uline{Remarque}. En théorie des groupes de BT, nous donnerons deux constructions de $\mathbb{D}^{*}(G)$, inspirées des deux constructions indiquées dans le cas des schémas abéliens: l'une ``cohomologique'', l'autre via une ``extension vectorielle universelle'' $\underline{E}(G)$. La construction de $\underline{E}(G)$ peut aussi s'obtenir par deux méthodes, l'une cohomologique, l'autre``par l'exponentielle'' étant directement adaptée à la démonstration du théorème de déformations.

